% =====================================================================
% Chapter 2 --- Theoretical Background                        [cross-cutting]
% Establishes ONCE the shared framework used by all later chapters, fixing
% the notation reused verbatim in Chapters 5--7. Two pillars: recursive
% Bayesian estimation (the tracker) and deep networks (the detector).
% =====================================================================

This chapter introduces the mathematical tools used throughout the dissertation
and fixes the notation reused in the later chapters. It rests on two pillars:
recursive Bayesian estimation, which underpins the tracking stage
(Chapters~\ref{cha:tracking} and \ref{cha:monocular}), and convolutional neural
networks, which underpin the detection front-end (Chapters~\ref{cha:pipeline}
and \ref{cha:airsim}). The treatment is deliberately self-contained but assumes
familiarity with linear algebra and probability.

\section{Bayesian Recursive Estimation and the Kalman Filter}

The tracking problem is an instance of \emph{recursive state estimation}: an
unobserved state evolves in time and is observed only through noisy
measurements. We model it by the discrete-time, possibly nonlinear, state-space
system
\begin{align}
  \bm{x}_k &= \bm{f}(\bm{x}_{k-1}) + \bm{w}_{k-1}, & \bm{w}_{k-1} &\sim \mathcal{N}(\bm{0},\bm{Q}_{k-1}), \label{eq:ss_process}\\
  \bm{z}_k &= \bm{h}(\bm{x}_k) + \bm{v}_k,         & \bm{v}_k     &\sim \mathcal{N}(\bm{0},\bm{R}_k),     \label{eq:ss_meas}
\end{align}
where $\bm{x}_k$ is the state, $\bm{z}_k$ the measurement, $\bm{f}(\cdot)$ the
process model, $\bm{h}(\cdot)$ the measurement model, and $\bm{w}_k,\bm{v}_k$
mutually independent zero-mean Gaussian process and measurement noises. The state
is assumed Markovian and the measurements conditionally independent given the
state.

Bayesian filtering propagates the posterior density $p(\bm{x}_k\mid \bm{z}_{1:k})$
in two alternating steps: a \emph{prediction} that pushes the previous posterior
through the process model (the Chapman--Kolmogorov equation), and an
\emph{update} that fuses the new measurement through Bayes' rule. For the
linear-Gaussian special case, with $\bm{f}(\bm{x})=\bm{F}\bm{x}$ and
$\bm{h}(\bm{x})=\bm{H}\bm{x}$, these integrals admit a closed form, the
\emph{Kalman filter} (KF) \cite{barshalom2001,brookner1998}: the posterior
remains Gaussian and is summarized by its mean $\hat{\bm{x}}_{k|k}$ and covariance
$\bm{P}_{k|k}$, propagated by
\begin{align}
  \hat{\bm{x}}_{k|k-1} &= \bm{F}\,\hat{\bm{x}}_{k-1|k-1}, &
  \bm{P}_{k|k-1} &= \bm{F}\,\bm{P}_{k-1|k-1}\,\bm{F}^{\top} + \bm{Q}, \label{eq:kf_predict}\\
  \bm{K}_k &= \bm{P}_{k|k-1}\bm{H}^{\top}\bm{S}_k^{-1}, &
  \bm{S}_k &= \bm{H}\,\bm{P}_{k|k-1}\,\bm{H}^{\top} + \bm{R}, \label{eq:kf_gain}\\
  \hat{\bm{x}}_{k|k} &= \hat{\bm{x}}_{k|k-1} + \bm{K}_k\bm{\nu}_k, &
  \bm{P}_{k|k} &= (\bm{I}-\bm{K}_k\bm{H})\,\bm{P}_{k|k-1}, \label{eq:kf_update}
\end{align}
with innovation $\bm{\nu}_k=\bm{z}_k-\bm{H}\hat{\bm{x}}_{k|k-1}$. The KF is the
optimal (minimum mean-square-error) estimator in the linear-Gaussian case.

\section{The Extended Kalman Filter}

The tracking models of this work are nonlinear, since the platform rotates and the
measurements are spherical or projective, so the linear KF does not apply
directly. The Extended Kalman Filter (EKF) retains the Gaussian parameterization
but linearizes the nonlinear models about the current estimate
\cite{terejanu2008,barshalom2001}.

\subsection{Linearization}

Let
\begin{equation}
  \bm{F}_k = \left.\frac{\partial \bm{f}(\bm{x})}{\partial \bm{x}}\right|_{\bm{x}=\hat{\bm{x}}_{k-1|k-1}},
  \qquad \mbox{and} \qquad
  \bm{H}_k = \left.\frac{\partial \bm{h}(\bm{x})}{\partial \bm{x}}\right|_{\bm{x}=\hat{\bm{x}}_{k|k-1}}
  \label{eq:ekf_jac}
\end{equation}
be the Jacobians of the process and measurement models evaluated at the predicted
state. The EKF propagates the covariance through these first-order
approximations.

\subsection{Prediction and Update}

The EKF's recursion mirrors \eqref{eq:kf_predict}--\eqref{eq:kf_update}, with the
nonlinear models used for the means and their Jacobians for the covariances:
\begin{align}
  \hat{\bm{x}}_{k|k-1} &= \bm{f}(\hat{\bm{x}}_{k-1|k-1}), &
  \bm{P}_{k|k-1} &= \bm{F}_k\,\bm{P}_{k-1|k-1}\,\bm{F}_k^{\top} + \bm{Q}, \label{eq:ekf_predict}\\
  \bm{\nu}_k &= \bm{z}_k - \bm{h}(\hat{\bm{x}}_{k|k-1}), &
  \bm{S}_k &= \bm{H}_k\,\bm{P}_{k|k-1}\,\bm{H}_k^{\top} + \bm{R}_k, \label{eq:ekf_innov}\\
  \bm{K}_k &= \bm{P}_{k|k-1}\bm{H}_k^{\top}\bm{S}_k^{-1}, &
  \hat{\bm{x}}_{k|k} &= \hat{\bm{x}}_{k|k-1} + \bm{K}_k\bm{\nu}_k, \label{eq:ekf_gain}
\end{align}
and $\bm{P}_{k|k} = (\bm{I}-\bm{K}_k\bm{H}_k)\,\bm{P}_{k|k-1}$. Unlike the linear
KF, the EKF is neither optimal nor unbiased: the linearization
\eqref{eq:ekf_jac} discards higher-order terms, and the magnitude of the
resulting error depends on how nonlinear the process and measurement models are
over the support of the state covariance. This caveat matters in Chapter~\ref{cha:tracking}, where two
algebraically distinct EKF formulations are compared.

\section{Coordinate Frames and Rigid-Body Transformations}

\subsection{World, Body, and Optical Frames}

Three right-handed frames are used. The \emph{world} frame $\mathcal{W}$ is a
local inertial East--North--Up (ENU) frame fixed to the operating area. The
\emph{body} frame $\mathcal{B}$ is attached to the platform (Forward--Left--Up).
The \emph{optical} frame $\mathcal{O}$ is attached to the camera (the computer
vision convention: $x$ right, $y$ down, $z$ forward). Estimation is carried in
$\mathcal{W}$ or $\mathcal{B}$ depending on the formulation, while measurements
originate in $\mathcal{O}$.

\subsection{Rotations, Euler Angles, and Homogeneous Transforms}

A rotation is an element of $SO(3)$, represented here by a $3\times3$ matrix
$\bm{R}$ with $\bm{R}^{\top}\bm{R}=\bm{I}$ and $\det\bm{R}=1$. Using the
$Z$--$Y$--$X$ (yaw--pitch--roll) Euler convention adopted throughout,
\begin{equation}
  \bm{R}(\phi,\theta,\psi) = \bm{R}_z(\psi)\,\bm{R}_y(\theta)\,\bm{R}_x(\phi),
  \label{eq:euler}
\end{equation}
with the elementary rotations
\begin{equation}
  \bm{R}_x(\phi)=\!\begin{bmatrix}1&0&0\\0&c_\phi&-s_\phi\\0&s_\phi&c_\phi\end{bmatrix}\!,\;
  \bm{R}_y(\theta)=\!\begin{bmatrix}c_\theta&0&s_\theta\\0&1&0\\-s_\theta&0&c_\theta\end{bmatrix}\!,\;
  \bm{R}_z(\psi)=\!\begin{bmatrix}c_\psi&-s_\psi&0\\s_\psi&c_\psi&0\\0&0&1\end{bmatrix}\!,
  \label{eq:elementary}
\end{equation}
where $c_\bullet=\cos$ and $s_\bullet=\sin$. With this convention $\bm{R}$ maps
body coordinates to world coordinates, and a point transforms as
$\bm{p}^{\mathcal{W}}=\bm{R}\,\bm{p}^{\mathcal{B}}+\bm{t}$; the inverse map is
$\bm{p}^{\mathcal{B}}=\bm{R}^{\top}(\bm{p}^{\mathcal{W}}-\bm{t})$. Equivalently, in
homogeneous coordinates,
\begin{equation}
  \begin{bmatrix} \bm{p}^{\mathcal{W}} \\ 1 \end{bmatrix}
  = \underbrace{\begin{bmatrix} \bm{R} & \bm{t} \\ \bm{0} & 1 \end{bmatrix}}_{\bm{T}\,\in\, SE(3)}
    \begin{bmatrix} \bm{p}^{\mathcal{B}} \\ 1 \end{bmatrix},
  \label{eq:se3_transform}
\end{equation}
which compactly writes the world point as a function of the body point through the
pose $\bm{T}\in SE(3)$.

\subsection{Pose Perturbations}

When the platform pose is uncertain (Chapter~\ref{cha:monocular}), its
uncertainty must be propagated through the geometry. A small pose perturbation is
represented in the tangent space of $SE(3)$ by a six-vector
$\bm{\xi}=[\bm{\delta\theta}^{\top}\;\bm{\delta t}^{\top}]^{\top}$, and to first
order a point in the camera frame perturbs as
\begin{equation}
  \delta\bm{p}^{\mathcal{O}} \approx -[\bm{p}^{\mathcal{O}}]_\times\,\bm{\delta\theta} + \bm{\delta t},
  \qquad
  [\bm{a}]_\times = \begin{bmatrix} 0 & -a_3 & a_2 \\ a_3 & 0 & -a_1 \\ -a_2 & a_1 & 0 \end{bmatrix},
  \label{eq:se3_pert}
\end{equation}
where $[\,\cdot\,]_\times$ is the skew-symmetric (cross-product) matrix
\cite{sola2018microlie}. This linearization is the basis of the extrinsic
Jacobian derived in Chapter~\ref{cha:monocular}.

\section{Motion Models}

\subsection{Constant-Velocity and Constant-Acceleration Models}

In the absence of a control input, target motion is described by a kinematic
model. The \emph{constant-velocity} (CV) model assumes piecewise-constant
velocity perturbed by random accelerations; per Cartesian axis, the
position--velocity pair $(p,\dot p)$ propagates as
\begin{equation}
  \begin{bmatrix} p_{k+1} \\ \dot p_{k+1} \end{bmatrix}
  = \begin{bmatrix} 1 & \Delta t \\ 0 & 1 \end{bmatrix}
    \begin{bmatrix} p_{k} \\ \dot p_{k} \end{bmatrix} + \text{noise},
  \label{eq:cv}
\end{equation}
with sampling interval $\Delta t$. The \emph{constant-acceleration} model adds an
explicit acceleration state. This work uses the CV model, with acceleration
represented as process noise rather than as a state variable.

\subsection{Process Noise}

Modelling the unknown acceleration as a continuous white noise of intensity
$\sigma_a^2$ yields, after discretization, the per-axis process-noise covariance
\begin{equation}
  \bm{Q}_{\text{axis}} = \sigma_a^{2}
  \begin{bmatrix} \Delta t^{3}/3 & \Delta t^{2}/2 \\[2pt] \Delta t^{2}/2 & \Delta t \end{bmatrix}.
  \label{eq:cwna}
\end{equation}
When the same intensity $\sigma_a$ is used on all axes, the resulting block-diagonal
$\bm{Q}$ is invariant under rotation of the coordinate axes, a property that has
a decisive consequence in Chapter~\ref{cha:tracking}.

\section{Measurement Models}

\subsection{Range--Bearing (Spherical) Measurements}

A depth or fused camera--LiDAR sensor reports the relative position of a target in
spherical coordinates. Writing the relative position in the sensor frame as
$\bm{p}^{c}=[x^{c},y^{c},z^{c}]^{\top}$, the measurement comprises range, azimuth
and elevation,
\begin{equation}
  \rho = \lVert\bm{p}^{c}\rVert,\quad
  \alpha = \operatorname{atan2}(y^{c},x^{c}),\quad
  \beta = \operatorname{atan2}\!\big(z^{c},\sqrt{(x^{c})^2+(y^{c})^2}\big).
  \label{eq:spherical}
\end{equation}
Ranging accuracy degrades with distance, which is captured by a heteroscedastic
range variance $\sigma_\rho^2 = \alpha_\rho\,\rho^2$, with $\alpha_\rho$ a
sensor-dependent coefficient \cite{fernandes2020}. This model is used by the
trackers of Chapters~\ref{cha:tracking}.

\subsection{Pinhole Projection and Lens Distortion}

A monocular camera (Chapter~\ref{cha:monocular}) instead measures the
\emph{pixel} projection of a world point. Under the pinhole model, a point
$\bm{X}^{\mathcal{O}}=[X,Y,Z]^{\top}$ in the optical frame projects to
\begin{equation}
  \bm{u} = \pi\!\big(\bm{K}\,\bm{X}^{\mathcal{O}}\big),
  \qquad
  \pi\!\left(\begin{bmatrix}X\\Y\\Z\end{bmatrix}\right) = \begin{bmatrix}X/Z\\Y/Z\end{bmatrix},
  \qquad
  \bm{K}=\begin{bmatrix} f_x & 0 & c_x \\ 0 & f_y & c_y \\ 0 & 0 & 1 \end{bmatrix},
  \label{eq:pinhole}
\end{equation}
where $\bm{K}$ is the intrinsic matrix ($f_x,f_y$ focal lengths, $(c_x,c_y)$
principal point). Real lenses introduce radial and tangential (plumb-bob)
distortion, parameterized by coefficients $(k_1,k_2,k_3,p_1,p_2)$, which must be
removed from each pixel measurement before it enters the filter; the procedure is
detailed in Chapter~\ref{cha:monocular}.

\section{Performance and Consistency Metrics}

Several families of metrics are used across the dissertation: the first three
evaluate the estimator, in terms of tracking accuracy, consistency, and
optimality, and the last evaluates the detection and tracking chain.

\subsection{Accuracy}
The root-mean-square error (RMSE) of the position estimate,
averaged over $M$ Monte~Carlo runs, summarizes accuracy:
\begin{equation}
  \mathrm{RMSE}(k)=\left(\frac1M\sum_{i=1}^{M}\big\lVert \bm{p}_k^{(i)}-\hat{\bm{p}}_k^{(i)}\big\rVert^2\right)^{1/2}.
  \label{eq:rmse}
\end{equation}

\subsection{Consistency}
A filter is \emph{consistent} when its reported
covariance matches its actual error. With ground truth available, the normalized
estimation error squared (NEES) of an $n$-dimensional state,
\begin{equation}
  \epsilon_k = (\bm{x}_k-\hat{\bm{x}}_{k|k})^{\top}\,\bm{P}_{k|k}^{-1}\,(\bm{x}_k-\hat{\bm{x}}_{k|k}),
  \label{eq:nees}
\end{equation}
follows a $\chi^2_n$ distribution (mean $n$) for a consistent filter. When ground
truth is unavailable, the normalized innovation squared (NIS),
$\eta_k=\bm{\nu}_k^{\top}\bm{S}_k^{-1}\bm{\nu}_k\sim\chi^2_m$ for $m$-dimensional
measurements, provides an analogous, measurement-only check \cite{barshalom2001}.

\subsection{Optimality Bound}
The posterior Cram\'er--Rao bound (PCRB) lower-bounds
the achievable estimation-error covariance. For the filtering problem it is
obtained from the Fisher information recursion
\begin{equation}
  \bm{J}_{k+1} = \big(\bm{Q}+\bm{F}\bm{J}_k^{-1}\bm{F}^{\top}\big)^{-1}
              + \mathbb{E}\!\big[\bm{H}_{k+1}^{\top}\bm{R}_{k+1}^{-1}\bm{H}_{k+1}\big],
  \label{eq:pcrb}
\end{equation}
where $\bm{J}_k$ is the Fisher information matrix. Its inverse lower-bounds the
estimation-error covariance of any unbiased estimator,
$\mathbb{E}[(\bm{x}_k-\hat{\bm{x}}_k)(\bm{x}_k-\hat{\bm{x}}_k)^{\top}]\succeq\bm{J}_k^{-1}$,
so the diagonal of $\bm{J}_k^{-1}$ bounds the per-component error variance from
below \cite{tichavsky1998}. This bound is used to assess near-optimality of the
tracker in Chapter~\ref{cha:tracking}.

\subsection{Detection and Multi-Object-Tracking Metrics}
\label{sec:detmetrics}

The detection front-end (Chapter~\ref{cha:airsim}) is evaluated with the standard
detection and tracking measures. A detection that matches a ground-truth object
under an overlap or distance gate is a true positive (TP) and otherwise a false
positive (FP); a ground-truth object with no matching detection is a false
negative (FN). \emph{Precision} ($P$), \emph{recall} ($R$), and their harmonic
mean $F_1$ are then
\begin{equation}
  P = \frac{\mathrm{TP}}{\mathrm{TP}+\mathrm{FP}}, \qquad
  R = \frac{\mathrm{TP}}{\mathrm{TP}+\mathrm{FN}}, \qquad
  F_1 = \frac{2\,P\,R}{P+R}.
  \label{eq:prf1}
\end{equation}
Sweeping the detector's confidence threshold traces a precision--recall curve
whose area is the average precision; averaging it over classes and overlap gates
gives the mean average precision (mAP), as standardized by the COCO benchmark
\cite{lin2014coco}.
For the temporal chain, multi-object-tracking accuracy (MOTA) aggregates false positives,
misses, and identity switches at one operating point, and the identity score
(IDF1) measures how consistently tracks preserve object identity. Because MOTA is
read at a single threshold, the driving community favors its integral over recall,
the average MOTA (AMOTA) \cite{weng2020ab3dmot}, which characterizes the
detector--tracker chain independently of a particular threshold; its companion
AMOTP averages the localization error of matched tracks. These measures are used
throughout Chapter~\ref{cha:airsim}.

\section{Deep Learning for Perception}

The detection front-end of the pipeline is built from neural networks; this
section summarizes the building blocks used in Chapters~\ref{cha:pipeline} and
\ref{cha:airsim}.

\subsection{Convolutional Neural Networks}

A convolutional neural network (CNN) applies learned convolutional filters across
an image, sharing weights spatially and stacking layers to build a hierarchy of
features, from edges and textures to object parts and whole objects. Weight
sharing makes CNNs parameter-efficient and translation-equivariant, which is why
they are the standard tool for image perception \cite{lecun2015deeplearning}.

\subsection{Two-Dimensional Object Detection and the YOLO Family}

Object detection localizes and classifies objects within an image. The decisive
split is between \emph{two-stage} detectors, such as the R-CNN family, which first
propose candidate regions and then classify and refine each one, and
\emph{one-stage} detectors, which predict all boxes and class scores in a single
forward pass. The You Only Look Once (YOLO) detector \cite{redmon2016yolo}
introduced the one-stage idea: the image is processed once by a fully
convolutional network that, over a dense spatial grid, directly regresses box
coordinates, an objectness score, and class probabilities. Removing the separate
proposal stage is what gives YOLO its real-time throughput, which is the property
that an onboard aerial system needs.

Architecturally, a modern YOLO network has three parts. A \emph{backbone} (a deep
convolutional feature extractor, in recent versions built from cross-stage-partial
blocks that split and merge feature maps to cut redundant computation) produces a
hierarchy of feature maps. A \emph{neck} (a feature-pyramid / path-aggregation
network) fuses these maps across scales so that both small and large objects are
represented. A \emph{detection head} then emits the dense predictions. Relative to
the original detector, the versions used today (the Ultralytics
YOLOv8--YOLOv11 line \cite{ultralytics2024yolo11}) introduce three changes that
matter for small aerial targets: they are \emph{anchor-free}, predicting box
offsets directly from grid locations instead of from hand-designed anchor boxes;
they use \emph{decoupled heads}, separating the classification and box-regression
branches; and they insert lightweight \emph{attention} modules in the backbone,
which reweight feature channels and spatial locations and improve the detection of
small, low-contrast objects, while preserving the single-pass design. A practical
difficulty in the aerial setting is \emph{domain shift}: because a detector's
accuracy depends on the match between its training distribution and the deployment
scene, a network meant to operate in an unseen environment is commonly trained on
a diverse union of datasets to widen its operating envelope.

\subsection{Point-Cloud Processing: PointNet and PointNet\texttt{++}}

Three-dimensional detection operates on point clouds, an \emph{unordered} set of
$N$ points $\{\bm{p}_i\in\mathbb{R}^3\}$ with no grid structure, so the convolution
that underpins image CNNs does not apply. PointNet \cite{qi2017pointnet} resolves
this with a construction that is invariant to the ordering of the points: each
point is passed independently through a \emph{shared} multilayer perceptron
$h(\cdot)$, and the resulting per-point features are aggregated by a
\emph{symmetric} function, max-pooling, into a single global descriptor
$\bm{g}=\max_i h(\bm{p}_i)$. Because max-pooling does not depend on the order of
its inputs, the network is permutation-invariant by construction; an optional
learned input alignment (a small ``T-Net'' that regresses a transformation
matrix) further canonicalizes the pose of the cloud. The limitation of this design
is that a single global pooling captures the overall shape but discards
\emph{local} neighborhood structure, which matters for fine geometry.

PointNet\texttt{++} \cite{qi2017pointnetpp} removes that limitation by applying
PointNet \emph{hierarchically}. Each \emph{set-abstraction} level (i) samples a
subset of centroid points by farthest-point sampling, (ii) groups the neighbors of
each centroid within a query radius (a ``ball query''), and (iii) encodes each
local group with a small PointNet. Stacking several such levels yields a
progressively coarser set of points carrying progressively richer local features,
so the network learns geometry at multiple scales, much as a CNN does on images.
This hierarchical point-set encoder is the workhorse of the 3D detector built in
this work. Frustum PointNets \cite{qi2018frustum} fuse it with the 2D detector:
each image detection defines a viewing frustum that crops the point cloud to a
small candidate region, and a PointNet\texttt{++} regresses the amodal 3D box
within that frustum, fusing the two modalities by geometry rather than in a learned
latent space.

\subsection{Single-Stage Voxel Detection}

The frustum approach is a \emph{two-stage} detector: a 2D network proposes
regions and a point-cloud network refines them in 3D, so its recall is bounded by
the 2D stage. The complementary \emph{single-stage} family detects directly in
the point cloud over the whole scene. Because raw points are irregular, these
methods first impose a regular structure: VoxelNet and the SECOND network
\cite{yan2018second} (a method name) quantize space into 3D voxels and apply (sparse) 3D
convolutions, while PointPillars \cite{lang2019pointpillars} collapses the
vertical axis into a bird's-eye-view (BEV) grid of \emph{pillars}, encodes each
pillar with a small PointNet, and runs an ordinary 2D convolutional backbone on
the resulting pseudo-image, trading a little vertical resolution for real-time
speed. Detection heads are increasingly \emph{center-based}: rather than
classifying anchor boxes, the network predicts a heatmap whose peaks are object
centers and regresses the box parameters at each peak \cite{zhou2019objects}.
This single-stage voxel paradigm is not part of the proposed pipeline; it is
included as the comparison baseline against which the two-stage frustum design is
evaluated on aerial data in Chapter~\ref{cha:airsim}.

\subsection{Training Losses for Detection}

Detectors are dominated by background: the vast majority of candidate locations
contain no object. Trained naively, the classifier is overwhelmed by easy
negatives. The \emph{focal loss} \cite{lin2017focal} addresses this by
down-weighting well-classified examples,
$\mathrm{FL}(p_t)=-\alpha_t\,(1-p_t)^{\gamma}\log p_t$, where $p_t$ is the
predicted probability of the true class and $\gamma>0$ concentrates the gradient
on hard examples. It is well suited to the objectness and heatmap classification
heads of dense detectors, whereas the continuous box parameters (center and size)
are regressed with a smooth-$L_1$ (Huber) loss.
